% GENERATED FILE. Edit canonical lessons, then run npm run generate:books.
% canonical-source-hash: fnv1a64:69a1156abd32c28e
% canonical-lessons: AR-C540-qiraa, AR-C540-ishtaraka, AR-C540-iltaqa, AR-C540-shaghaf, AR-C540-wala

\chapter{Reading, clubs and passions}
\label{ch:ar-reading-clubs-and-passions}
\hlchaptermodality{540}

\begin{chapteropening}
\textbf{By the end of this chapter:} \emph{I can talk about reading, passion and fondness, and say that someone subscribed or met up with a friend.}

Complete the last lesson of chapter 540: I can talk about reading, passion and fondness, and say that someone subscribed or met up with a friend.
\end{chapteropening}

\section[\texorpdfstring{\ar{قراءة} (qirāʾa)}{قراءة (qirāʾa)}]{\ar{قراءة} (qirāʾa) — reading}

\label{lesson:AR-C540-qiraa}



\begin{quote}
Before the new one: say the Arabic for gardening, then the Arabic for to swing.
\end{quote}

\subsection*{You'll want to know: \ar{قراءة}}
\textbf{\ar{قراءة}} — \emph{qirāʾa} — \textquotedblleft{}reading\textquotedblright{}.

\begin{quote}
\ar{القراءة} \ar{هوايتي} \ar{المفضلة}.
\end{quote}

\begin{quote}
\emph{al-qirāʾa hiwāyatī al-mufaḍḍala.}
\end{quote}

\begin{quote}
Reading is my favourite hobby.
\end{quote}

\begin{cousinweb}
\textbf{\ar{قراءة}} (\emph{qirāʾa}) is built on the native Arabic root \emph{q-r-a}, \textquotedblleft{}to read, to recite\textquotedblright{} (Lane). Qirāʾa is the verbal noun of qaraʾa, to read, which you know.
\end{cousinweb}

\begin{grammarlens}[title={What you've built}]
One more word for describing what you lived and felt.
\end{grammarlens}

\subsection*{Your turn}
\begin{itemize}
\raggedright
  \item \emph{Say it:} \emph{qirāʾa}
  \item \emph{Say it:} \emph{qirāʾa}, once more
  \item \emph{Say it:} \emph{al-qirāʾa hiwāyatī al-mufaḍḍala.} — Reading is my favourite hobby.
  \item \emph{Recall:} say the Arabic for gardening, then the Arabic for to swing, then say \emph{qirāʾa} again
\end{itemize}

\begin{tcolorbox}[breakable,colback=teal!4,colframe=teal!35!black,title={Before you move on}]
What does \ar{قراءة} mean? (\textquotedblleft{}Reading\textquotedblright{}.) Say it once more.
\end{tcolorbox}
\section[\texorpdfstring{\ar{اشترك} (ishtaraka)}{اشترك (ishtaraka)}]{\ar{اشترك} (ishtaraka) — to subscribe; to take part}

\label{lesson:AR-C540-ishtaraka}



\begin{quote}
Before the new one: say the Arabic for to swing, then the Arabic for reading.
\end{quote}

\subsection*{You'll want to know: \ar{اشترك}}
\textbf{\ar{اشترك}} — \emph{ishtaraka} — \textquotedblleft{}to subscribe; to take part\textquotedblright{}.

\begin{quote}
\ar{اشترك} \ar{أخي} \ar{في} \ar{مجلة} \ar{رياضية}.
\end{quote}

\begin{quote}
\emph{ishtaraka akhī fī majalla riyāḍiyya.}
\end{quote}

\begin{quote}
My brother subscribed to a sports magazine.
\end{quote}

Present: \ar{يشترك}, \emph{yashtarik}.

It takes \ar{في}.

\begin{cousinweb}
\textbf{\ar{اشترك}} (\emph{ishtaraka}) is built on the native Arabic root \emph{sh-r-k}, \textquotedblleft{}to share, to be a partner\textquotedblright{} (Lane). Sharika, a company, and shāraka, to share, which you know, are from the same root.
\end{cousinweb}

\begin{grammarlens}[title={What you've built}]
One more word for describing what you lived and felt.
\end{grammarlens}

\subsection*{Your turn}
\begin{itemize}
\raggedright
  \item \emph{Say it:} \emph{ishtaraka}
  \item \emph{Say it:} \emph{ishtaraka}, once more
  \item \emph{Say it:} \emph{ishtaraka akhī fī majalla riyāḍiyya.} — My brother subscribed to a sports magazine.
  \item \emph{Recall:} say the Arabic for to swing, then the Arabic for reading, then say \emph{ishtaraka} again
\end{itemize}

\begin{tcolorbox}[breakable,colback=teal!4,colframe=teal!35!black,title={Before you move on}]
What does \ar{اشترك} mean? (\textquotedblleft{}To subscribe; to take part\textquotedblright{}.) Say it once more.
\end{tcolorbox}
\section[\texorpdfstring{\ar{التقى} (iltaqā)}{التقى (iltaqā)}]{\ar{التقى} (iltaqā) — to meet up (with)}

\label{lesson:AR-C540-iltaqa}



\begin{quote}
Before the new one: say the Arabic for reading, then the Arabic for to subscribe.
\end{quote}

\subsection*{You'll want to know: \ar{التقى}}
\textbf{\ar{التقى}} — \emph{iltaqā} — \textquotedblleft{}to meet up (with)\textquotedblright{}.

\begin{quote}
\ar{التقيت} \ar{بزميلي} \ar{في} \ar{المقهى}.
\end{quote}

\begin{quote}
\emph{iltaqaytu bi-zamīlī fī al-maqhā.}
\end{quote}

\begin{quote}
I met up with my colleague at the café.
\end{quote}

Present: \ar{يلتقي}, \emph{yaltaqī}.

It takes \ar{بـ}, or a direct object.

\begin{cousinweb}
\textbf{\ar{التقى}} (\emph{iltaqā}) is built on the native Arabic root \emph{l-q-y}, \textquotedblleft{}to meet; to throw\textquotedblright{} (Lane). Ilā l-liqāʾ, see you, which you know, is from the same root.
\end{cousinweb}

\begin{grammarlens}[title={What you've built}]
One more word for describing what you lived and felt.
\end{grammarlens}

\subsection*{Your turn}
\begin{itemize}
\raggedright
  \item \emph{Say it:} \emph{iltaqā}
  \item \emph{Say it:} \emph{iltaqā}, once more
  \item \emph{Say it:} \emph{iltaqaytu bi-zamīlī fī al-maqhā.} — I met up with my colleague at the café.
  \item \emph{Recall:} say the Arabic for reading, then the Arabic for to subscribe, then say \emph{iltaqā} again
\end{itemize}

\begin{tcolorbox}[breakable,colback=teal!4,colframe=teal!35!black,title={Before you move on}]
What does \ar{التقى} mean? (\textquotedblleft{}To meet up (with)\textquotedblright{}.) Say it once more.
\end{tcolorbox}
\section[\texorpdfstring{\ar{شغف} (shaghaf)}{شغف (shaghaf)}]{\ar{شغف} (shaghaf) — passion}

\label{lesson:AR-C540-shaghaf}



\begin{quote}
Before the new one: say the Arabic for to subscribe, then the Arabic for to meet up.
\end{quote}

\subsection*{You'll want to know: \ar{شغف}}
\textbf{\ar{شغف}} — \emph{shaghaf} — \textquotedblleft{}passion\textquotedblright{}.

\begin{quote}
\ar{يتكلم} \ar{المعلم} \ar{عن} \ar{الأدب} \ar{بشغف}.
\end{quote}

\begin{quote}
\emph{yatakallamu al-muʿallim ʿan al-adab bi-shaghaf.}
\end{quote}

\begin{quote}
The teacher talks about literature with passion.
\end{quote}

\begin{cousinweb}
\textbf{\ar{شغف}} (\emph{shaghaf}) is built on the native Arabic root \emph{sh-gh-f}, \textquotedblleft{}the heart's covering; to strike the heart\textquotedblright{} (Lane). Lane gives shaghaf as love that reaches the shaghāf, the covering of the heart.
\end{cousinweb}

\begin{grammarlens}[title={What you've built}]
One more word for describing what you lived and felt.
\end{grammarlens}

\subsection*{Your turn}
\begin{itemize}
\raggedright
  \item \emph{Say it:} \emph{shaghaf}
  \item \emph{Say it:} \emph{shaghaf}, once more
  \item \emph{Say it:} \emph{yatakallamu al-muʿallim ʿan al-adab bi-shaghaf.} — The teacher talks about literature with passion.
  \item \emph{Recall:} say the Arabic for to subscribe, then the Arabic for to meet up, then say \emph{shaghaf} again
\end{itemize}

\begin{tcolorbox}[breakable,colback=teal!4,colframe=teal!35!black,title={Before you move on}]
What does \ar{شغف} mean? (\textquotedblleft{}Passion\textquotedblright{}.) Say it once more.
\end{tcolorbox}
\section[\texorpdfstring{\ar{ولع} (walaʿ)}{ولع (walaʿ)}]{\ar{ولع} (walaʿ) — fondness, keen love}

\label{lesson:AR-C540-wala}



\begin{quote}
Before the new one: say the Arabic for to meet up, then the Arabic for passion.
\end{quote}

\subsection*{You'll want to know: \ar{ولع}}
\textbf{\ar{ولع}} — \emph{walaʿ} — \textquotedblleft{}fondness, keen love\textquotedblright{}.

\begin{quote}
\ar{عند} \ar{ابني} \ar{ولع} \ar{بالسيارات}.
\end{quote}

\begin{quote}
\emph{ʿinda ibnī walaʿ bi-s-sayyārāt.}
\end{quote}

\begin{quote}
My son is crazy about cars.
\end{quote}

\ar{مولع} \ar{بـ} is keen on.

\begin{cousinweb}
\textbf{\ar{ولع}} (\emph{walaʿ}) is built on the native Arabic root \emph{w-l-ʿ}, \textquotedblleft{}to make eagerly desirous\textquotedblright{} (Lane). Lane gives Form IV awlaʿa as to make someone eagerly fond of a thing.
\end{cousinweb}

\begin{grammarlens}[title={What you've built}]
One more word for describing what you lived and felt.
\end{grammarlens}

\subsection*{Your turn}
\begin{itemize}
\raggedright
  \item \emph{Say it:} \emph{walaʿ}
  \item \emph{Say it:} \emph{walaʿ}, once more
  \item \emph{Say it:} \emph{ʿinda ibnī walaʿ bi-s-sayyārāt.} — My son is crazy about cars.
  \item \emph{Recall:} say the Arabic for to meet up, then the Arabic for passion, then say \emph{walaʿ} again
\end{itemize}

\begin{tcolorbox}[breakable,colback=teal!4,colframe=teal!35!black,title={Before you move on}]
What does \ar{ولع} mean? (\textquotedblleft{}Fondness, keen love\textquotedblright{}.) Say it once more.
\end{tcolorbox}
